\section{A numerical threshold for the two independent-copy tails}
\label{sec:iid_tail_threshold}

The following arithmetic estimate supplies the eventual threshold used when
combining the two tail bounds in \cite[Section~7]{OpenAI2026AreaLaw}.
The coefficient and polynomial degree are fixed as the number of copies
increases.

\begin{theorem}[An eventual error bound]
\label{thm:iid_tail_threshold}
\lean{Real.eventually_iid_tail_error_le_half}
\leanok
For every real number $V$ and every integer $m\ge0$, there is an integer
$K\ge1$ such that, for all integers $k\ge K$,
\begin{align}
 \frac{V}{\sqrt{k}}+(k+1)^m\exp\!\left(-\frac{k^{3/4}}{2}\right)
   &\le \frac12.
 \label{eq:iid_tail_threshold}
\end{align}
In particular, the estimate applies to a nonnegative variance $V$.
It requires no spectral or representation label assumptions.
\end{theorem}
\begin{proof}
\leanok
For $k\ge1$, $(k+1)^m\le 2^mk^m$. Put $t=k^{3/4}$, so that
$k^m=t^{4m/3}$. The standard exponential decay estimate gives
\begin{align}
 t^{4m/3}e^{-t/2}&\longrightarrow0
 \qquad(t\longrightarrow\infty).
 \label{eq:iid_tail_threshold_decay}
\end{align}
Consequently, $2^mk^me^{-k^{3/4}/2}$ tends to zero as $k$ tends to
infinity. The term $V/\sqrt{k}$ also tends to zero. Their sum is therefore
eventually less than $1/2$; domination by this sum proves the displayed
inequality, and increasing $K$ ensures $K\ge1$.
\end{proof}
